\subsection{The uniform SAT and exact-counting algorithm}
\label{sec:counting-algorithm}

We now prove the algorithmic guarantee highlighted in
Section~\ref{sec:algorithm-overview}. The nonuniform theorem used
synthesis on formal ordinary inputs.
The all-quantified case has no such step. Its explicit table operations
also preserve counts, which gives the following algorithmic consequence.
This statement includes its proof; the unresolved issue in
Section~\ref{sec:prior-work} is comparison and priority.

\begin{corollary}[Counting by the gate-minus-input budget]
\label{cor:gate-counting}
Fix $\delta>0$ and write $\alpha=1/3+\delta$. Given a single-output
$B_2$ circuit with $s$ gates and $u$ input bits, one can count its
satisfying inputs in time
\[
 \poly_\delta(s+u+1)\,
 2^{\min\{b,\alpha(s+1-b)\}},
\]
where $b$ is the number of inputs remaining in a normalized output
cone. In particular, for every fixed $\varepsilon>0$, the time is
\[
 \poly_\varepsilon(s+u+1)\,
 2^{(1/4+\varepsilon)(s+1)}.
\]
The space used is at most polynomial times the same exponential
factor. A polynomial-space bound is not asserted.
\end{corollary}
\begin{proof}
Prune the output cone and propagate constants and unary redundancies
without changing the function. Inputs removed from the cone are free;
their contribution is a final multiplicative factor $2^{u-b}$.
Constant and designated-input outputs are evaluated directly. For the
remaining circuit quantify all $b$ inputs, so there are no independent
signals. Its consistency rank satisfies $\kappa\le s+1-b$.

Replace Boolean OR and AND in the tensor construction by integer
addition and multiplication. The original tables contain only zeros
and ones. For each satisfying input, the circuit equations specify
exactly one gate-value assignment. Incidence equality tables specify
exactly one set of equal edge values, and their ternary expansions have
exactly one internal extension. Thus the initial tensor contraction
counts satisfying inputs, with no spurious multiplicity.

Every reduction in Lemma~\ref{lem:cubic-reduction} is a distributive
identity, including diagonal loop sums and joint parallel-edge sums.
It remains valid over the nonnegative integers. More explicitly, a
reduced table entry is the number of internal edge assignments in its
contracted region that satisfy the original tables, conditional on its
boundary. A contraction combines disjoint regions and sums over their
shared indices, preserving this interpretation. The same invariant
holds for the frontier entries in Lemma~\ref{lem:frontier}.

There are $O(s+1)$ indices after equality expansion. Each table entry
is therefore bounded by $2^{O(s+1)}$, so its bit length is $O(s+1)$.
Integer additions and multiplications have polynomial bit cost; the
last factor $2^{u-b}$ only shifts the answer. For fixed $\delta$, the
cited constructive pathwidth theorem and median conversion find the
required order in polynomial time. The bounded-size graphs below its
threshold are handled by exhaustive order search, a constant depending
on $\delta$. The operation count is consequently polynomial times
$2^{\alpha(s+1-b)}$, with exponential-size frontiers allowed.

Alternatively enumerate the $2^b$ used-input assignments, evaluate the
circuit, and count the accepting ones. Choose the cheaper proved bound.
Finally,
\[
 \min\{b,\alpha(s+1-b)\}
 \le\frac{\alpha}{1+\alpha}(s+1),
 \qquad \frac{\alpha}{1+\alpha}\longrightarrow\frac14.
\]
This proves both time estimates and the stated space majorant.
\end{proof}

Testing whether the count is nonzero decides circuit satisfiability.
For any fixed $0<\gamma<4$, the bound also gives a running time
$2^{(1-\Omega_\gamma(1))u}\poly(u)$ when
$s\le(4-\gamma)u$, after choosing a sufficiently small fixed
$\varepsilon$. This consequence uses exponential space and the exact
$B_2$ gate convention. Its numerical comparison with previously
reported circuit-SAT algorithms deserves particular scrutiny.
